\uuid{mK9p}
\titre{Tests d'hypothèses et test d'adéquation}
\niveau{L3}
\module{Probabilité et statistique}
\chapitre{Tests sur la moyenne, test du chi2}
\sousChapitre{Tests d'hypothèses, intervalle de confiance}
\theme{Test de Student, puissance, test du chi-deux, loi binomiale}
\auteur{Grégoire Menet}
\datecreate{2026-05-19}
\organisation{AMSCC}
\difficulte{3}

\contenu{
	
	\texte{
		Emmanuel, étudiant en ESM3, se demande combien de temps ses camarades ont révisé en moyenne pour l'examen de statistique inférentielle. Il estime que ce temps moyen est d'environ 3 heures. Mathis n'est pas d'accord : selon lui, les étudiants ont révisé davantage qu'ils ne le laissent paraître.
		
		Pour départager leurs points de vue, Emmanuel interroge 6 camarades et obtient les temps de révision suivants (en heures) :
		$$
		0{,}5 \quad 5 \quad 1 \quad 8 \quad 3 \quad 0{,}083
		$$
		On suppose que le temps de révision suit une loi normale $\mathcal{N}(\mu, \sigma)$ avec $\mu$ et $\sigma$ inconnus.
	}
	
	\begin{enumerate}
		
		\item
		\question{Construire un test statistique permettant de départager les hypothèses d'Emmanuel ($\mu = 3$ h) et de Mathis ($\mu > 3$ h), au risque de première espèce $\alpha = 5\%$.}
		\indication{La variance étant inconnue, on utilisera la statistique de Student.}
		\reponse{
			On teste $H_0 : \mu = 3$ contre $H_1 : \mu > 3$. Il s'agit d'un test unilatéral à droite.
			
			On calcule la moyenne empirique et l'écart-type corrigé sur l'échantillon :
			\[
			\bar{x} = \frac{0{,}5 + 5 + 1 + 8 + 3 + 0{,}083}{6} \approx 2{,}93 \text{ h},
			\qquad
			s^2 = \frac{1}{n-1}\sum_{i=1}^{6}(x_i - \bar{x})^2 \approx 8{,}34.
			\]
			La statistique de test est :
			\[
			T = \frac{\bar{X} - \mu_0}{S/\sqrt{n}} \sim \mathcal{S}t(n-1) = \mathcal{S}t(5) \quad \text{sous } H_0.
			\]
			La région de rejet est $\mathcal{R} = \left]q_{5;\,1-\alpha} ; +\infty\right[$ où $q_{5;\,0{,}95} \approx 2{,}015$.
			
			Application numérique :
			\[
			t_{\text{obs}} = \frac{2{,}93 - 3}{\sqrt{8{,}34}/\sqrt{6}} \approx -0{,}059.
			\]
			Comme $t_{\text{obs}} \approx -0{,}059 < 2{,}015$, on ne rejette pas $H_0$ au risque $5\%$ : les données ne permettent pas de conclure que les étudiants ont révisé davantage que 3 heures.
		}
		
		\item
		\question{Mathis pense désormais que le temps moyen de révision est plutôt de $\mu_1 = 4{,}5$ h. En supposant la variance connue et égale à $\sigma^2 = 8$ h$^2$, calculer la puissance du test $H_0 : \mu = 3$ contre $H_1 : \mu = 4{,}5$ au niveau $\alpha = 5\%$. Commenter.}
		\indication{La puissance est la probabilité de rejeter $H_0$ lorsque $\mu = \mu_1$.}
		\reponse{
			La variance étant maintenant connue, la statistique de test est :
			\[
			Z = \frac{\bar{X} - 3}{\sigma/\sqrt{n}} = \frac{\bar{X} - 3}{\sqrt{8}/\sqrt{6}} \sim \mathcal{N}(0,1) \quad \text{sous } H_0.
			\]
			La région de rejet est $\mathcal{R} = \left]z_{0{,}95} ; +\infty\right[$ avec $z_{0{,}95} \approx 1{,}645$.
			
			La puissance est la probabilité de rejeter $H_0$ lorsque $\mu = 4{,}5$ :
			\begin{align*}
				\pi(\mu_1)
				&= P_{\mu_1}\!\left(Z > 1{,}645\right)
				= P_{\mu_1}\!\left(\frac{\bar{X} - 3}{\sqrt{8/6}} > 1{,}645\right) \\
				&= P\!\left(\mathcal{N}(0,1) > 1{,}645 - \frac{4{,}5 - 3}{\sqrt{8/6}}\right)
				= P\!\left(\mathcal{N}(0,1) > 1{,}645 - 1{,}837\right) \\
				&= P\!\left(\mathcal{N}(0,1) > -0{,}192\right)
				\approx \Phi(0{,}192) \approx 0{,}576.
			\end{align*}
			La puissance est d'environ $57{,}6\%$ : avec seulement $n = 6$ observations, le test a moins d'une chance sur deux de détecter un écart de $1{,}5$ heure par rapport à $H_0$. L'échantillon est trop petit pour un test puissant.
		}
		
		\item
		\texte{
			Après l'examen, Emmanuel et Mathis s'interrogent sur la distribution des notes. Ils souhaitent tester l'hypothèse selon laquelle les notes suivent une loi binomiale $\mathcal{B}\!\left(5;\frac{1}{2}\right)$. L'examen est noté sur 5 et 96 étudiants ont composé. Le tableau des effectifs observés est le suivant :
			
			\begin{center}
				\begin{tabular}{|c|c|c|c|c|c|c|}
					\hline
					Note & 0 & 1 & 2 & 3 & 4 & 5 \\
					\hline
					Effectif observé & 1 & 5 & 33 & 44 & 12 & 1 \\
					\hline
				\end{tabular}
			\end{center}
		}
		\question{À l'aide d'un test du $\chi^2$ d'adéquation au risque $\alpha = 2{,}5\%$, tester l'hypothèse que les notes suivent une loi $\mathcal{B}\!\left(5;\frac{1}{2}\right)$. On pensera à regrouper les notes $\{0,1\}$ et $\{4,5\}$ pour éviter les effectifs trop faibles dans une classe.}
		\indication{Calculer les effectifs théoriques sous $H_0$, vérifier qu'ils sont tous supérieurs à 5 (regrouper si nécessaire), puis calculer la statistique $\chi^2$.}
		\reponse{
			Sous $H_0 : X \sim \mathcal{B}(5;\frac{1}{2})$, les probabilités théoriques sont :
			\[
			p_k = \binom{5}{k}\left(\frac{1}{2}\right)^5, \quad k = 0,\ldots,5.
			\]
			Les effectifs théoriques $e_k = 96 \times p_k$ sont :
			
			\begin{center}
				\begin{tabular}{|c|c|c|c|c|c|c|}
					\hline
					Note $k$ & 0 & 1 & 2 & 3 & 4 & 5 \\
					\hline
					$p_k$ & $1/32$ & $5/32$ & $10/32$ & $10/32$ & $5/32$ & $1/32$ \\
					\hline
					$e_k$ & $3{,}00$ & $15{,}00$ & $30{,}00$ & $30{,}00$ & $15{,}00$ & $3{,}00$ \\
					\hline
				\end{tabular}
			\end{center}
			
			Les effectifs théoriques des notes 0 et 5 sont inférieurs à 5 ; on les regroupe avec leurs voisins :
			
			\begin{center}
				\begin{tabular}{|c|c|c|c|c|}
					\hline
					Classe & $\{0,1\}$ & $\{2\}$ & $\{3\}$ & $\{4,5\}$ \\
					\hline
					Effectif observé $o_j$ & 6 & 33 & 44 & 13 \\
					\hline
					Effectif théorique $e_j$ & $18{,}00$ & $30{,}00$ & $30{,}00$ & $18{,}00$ \\
					\hline
				\end{tabular}
			\end{center}
			
			La statistique de test est :
			\[
			D = \sum_{j=1}^{4} \frac{(o_j - e_j)^2}{e_j}
			= \frac{(6-18)^2}{18} + \frac{(33-30)^2}{30} + \frac{(44-30)^2}{30} + \frac{(13-18)^2}{18}
			= 8{,}00 + 0{,}30 + 6{,}53 + 1{,}39 \approx 16{,}22.
			\]
			Sous $H_0$, $D \sim \chi^2(r - 1)$ où $r = 4$ est le nombre de classes après regroupement et les paramètres de la loi sont supposés connus (non estimés). Donc $D \sim \chi^2(3)$.
			
			Le quantile $\chi^2_{3;\,0{,}975} \approx 9{,}35$. Comme $D \approx 16{,}22 > 9{,}35$, on rejette $H_0$ au risque $2{,}5\%$ : les données ne sont pas compatibles avec une loi $\mathcal{B}(5;\frac{1}{2})$.
		}
		
	\end{enumerate}
}